% Part: incompleteness
% Chapter: representability-in-q
% Section: beta-function

\documentclass[../../../include/open-logic-section]{subfiles}

\begin{document}

\olfileid{inc}{req}{bet}
\olsection{બીટા વિધેયનું લેમ્મા}


સરવાળો, ગુણાકાર અને $\Char{=}$
ઉપલબ્ધ હોય ત્યારે આદિમ પુનરાવર્તન
કરી શકાય છે તે બતાવવા આપણે શ્રેણીઓ
સંભાળતાં વિધેયો વિકસાવવાં પડશે.
(ઘાતાંકન પણ ઉપલબ્ધ હોત, તો આ કામ
સરળ બનત.) આદિમ પુનરાવર્તન
ઉપલબ્ધ હતું ત્યારે ``$n$મો
અવિભાજ્ય'' જેવા વિધેયો વ્યાખ્યાયિત
કરીને પ્રમાણમાં સીધું સંકેતીકરણ
પસંદ કરી શકતા હતા. પરંતુ અહીં
આદિમ પુનરાવર્તન ઉપલબ્ધ નથી—હકીકતમાં
લઘુત્તમીકરણ વડે આદિમ પુનરાવર્તન
કરી શકાય છે તે જ બતાવવાનું છે—તેથી
વધુ ચતુર રીત જોઈએ.

\begin{lem}
\ollabel{lem:beta}
એવું વિધેય $\beta(d,i)$ છે કે
દરેક શ્રેણી $a_0$, \dots,~$a_n$
માટે એવી સંખ્યા~$d$ મળે છે કે
દરેક $i \le n$ માટે
$\beta(d,i) = a_i$. વળી, $\beta$
મૂળભૂત વિધેયોમાંથી માત્ર સંયોજન
અને નિયમિત લઘુત્તમીકરણ વડે
વ્યાખ્યાયિત કરી શકાય છે.
\end{lem}

$d$ને શ્રેણી $\tuple{a_0, \dots,
a_n}$નો સંકેત અને $\beta(d,i)$ને
તેનો $i$મો ઘટક પરત આપતું વિધેય
સમજો. (ધ્યાન રહે કે આ
``સંકેતીકરણ'' આપણે અગાઉ જાણેલા
અવિભાજ્ય સંખ્યાઓની ઘાતોવાળા સંકેતીકરણનો
ઉપયોગ \emph{કરતું નથી}!) લેમ્મા
ખૂબ મર્યાદિત દાવો કરે છે: શ્રેણીઓ
જોડી શકાય કે ઘટકો ઉમેરી શકાય
એવું તે કહેતું નથી. સંયોજન અને
નિયમિત લઘુત્તમીકરણ વડે વ્યાખ્યાયિત
વિધેયોનો ઉપયોગ કરીને $a_0$,
\dots,~$a_n$માંથી~$d$ \emph{ગણી}
શકાય એવું પણ તે કહેતું નથી.
તે માત્ર એટલું કહે છે કે એવું
``સંકેત ઉકેલતું'' વિધેય છે જેના માટે
દરેક શ્રેણી ``સંકેતિત'' થાય છે.

$\beta$ સંકેતનો ઉપયોગ ગડેલે કર્યો
હતો. ફરી કહીએ તો, લેમ્માની
સાબિતીમાં અઘરું કામ દેખીતી રીતે
મર્યાદિત સાધનો—માત્ર સંયોજન અને
લઘુત્તમીકરણ—વડે યોગ્ય~$\beta$
વ્યાખ્યાયિત કરવાનું છે. જોકે સરવાળો,
ગુણાકાર અને~$\Char{=}$ વાપરી
શકીએ છીએ. આ લેમ્માને સાબિત
કરવાની અનેક રીતો છે; તેમાંની એક
સૌથી સરળ રીત ગડેલની મૂળ રીત છે.
તેમાં સુનઝીનું પ્રમેય કહેવાતું
સંખ્યાસિદ્ધાંતનું તથ્ય વપરાય છે
(પરંપરાગત નામ ``ચીની શેષફળ પ્રમેય'').

\begin{defn}
બે પ્રાકૃતિક સંખ્યાઓ $a$ અને $b$
તેમનો મહત્તમ સામાન્ય અવયવ~$1$
હોય ત્યારે અને માત્ર ત્યારે
\emph{પરસ્પર અવિભાજ્ય} કહેવાય છે;
બીજા શબ્દોમાં, ૧ સિવાયનો
કોઈ અવયવ બંનેમાં સામાન્ય નથી.
\end{defn}

\begin{defn}
પ્રાકૃતિક સંખ્યાઓ $a$ અને $b$
\emph{મોડ્યુલો~$c$ પ્રમાણે સમશેષ}
છે, $a \equiv b \mod c$, ત્યારે
અને માત્ર ત્યારે $c \mid (a-b)$;
એટલે કે~$c$ વડે ભાગતાં
$a$ અને~$b$નું
શેષફળ સમાન મળે છે.
\end{defn}

સુનઝીનું પ્રમેય આ છે:
\begin{thm}
માનો કે $x_0$, \dots,~$x_n$
કોઈપણ બે પરસ્પર અવિભાજ્ય છે.
$y_0$, \dots,~$y_n$ કોઈપણ
સંખ્યાઓ લો. તો એવી સંખ્યા $z$
છે કે
\begin{align*}
z & \equiv y_0 \mod x_0 \\
z & \equiv y_1 \mod x_1 \\
& \vdots  \\
z & \equiv y_n \mod x_n.
\end{align*}
\end{thm}

સુનઝીના પ્રમેયનો ઉપયોગ આમ કરીશું:
જો $x_0$, \dots,~$x_n$ અનુક્રમે
$y_0$, \dots,~$y_n$ કરતાં મોટાં
હોય, તો $z$ને શ્રેણી
$\tuple{y_0, \dots,y_n}$નો
સંકેત લઈ શકાય. $y_i$ પાછું
મેળવવા $z$ને~$x_i$ વડે ભાગીને
શેષફળ લેવું પૂરતું છે. આ
સંકેતીકરણ માટે $x_0$,
\dots,~$x_n$નાં યોગ્ય મૂલ્યો
શોધવાં પડશે.

બે અવલોકનો તેમાં મદદ કરશે.
$y_0$, \dots,~$y_n$ આપેલાં હોય,
તો આ લો:
\begin{align*}
j &= \max(n, y_0 + 1, \dots, y_n + 1), \\
m &= \lcm(1,\dots,j),
\end{align*}
અને પછી આ લો:
\begin{align*}
x_0 & = 1 + m \\
x_1 & = 1 + 2 \cdot m \\
x_2 & = 1 + 3 \cdot m \\
& \vdots  \\
x_n & = 1 + (n+1) \cdot m
\end{align*}
તો બે બાબતો સાચી છે:
\begin{enumerate}
\item\ollabel{rel-prime} $x_0,\dots,x_n$માંથી કોઈપણ બે પરસ્પર અવિભાજ્ય છે.
\item\ollabel{less} દરેક $i$ માટે $y_i < x_i$.
\end{enumerate}
\olref{rel-prime} કેમ સાચું છે તે
જોવા માનો કે $i\ne k$, $p$ અવિભાજ્ય
છે અને
$p \mid x_i$ તથા $p \mid x_k$.
\sourcecorrection{OLINC-024}{સ્થિર અંગ્રેજી
મૂળની આ દલીલમાં $i$ અને~$k$
જુદા છે એમ સ્પષ્ટ કરવું જરૂરી
છે; નહીં તો $i-k=0$ પરથી
$p\mid m$ નીકળતું નથી.
પરસ્પર અવિભાજ્યતાના દાવા
મુજબ અહીં $i\ne k$ લીધું છે.}
તો $p \mid 1 + (i+1) m$ અને
$p \mid 1 + (k+1) m$; તેથી $p$
તેમના તફાવતનો પણ અવયવ છે:
\[
(1 + (i+1)m) - (1+ (k+1)m) = (i-k) m.
\]
$p$ એ $1 + (i+1)m$નો અવયવ
હોવાથી તે~$m$નો અવયવ હોઈ
શકે નહીં (નહીં તો પહેલું પદ~$p$
વડે ભાગતાં શેષ~$1$ રહે).
હવે $p$ એ $(i-k)m$નો અવયવ
છે અને ઉપર પ્રમાણે બીજા
ગુણકનો નથી, તેથી તે
$i-k$નો અવયવ છે. પરંતુ
$\left|i-k\right|$ વધુમાં વધુ~$n$
છે અને $j \geq n$ પસંદ કર્યું છે;
એટલે $p \mid m$, ફરી વિરોધાભાસ.
આથી $x_i$ અને $x_k$ બંનેનો
અવયવ હોય એવી અવિભાજ્ય સંખ્યા નથી.
\olref{less} સરળ છે:
$y_i < j \leq m < x_i$.

હવે $\beta$-વિધેયનું લેમ્મા સાબિત
કરીએ. યાદ રાખો કે $0$, ઉત્તરગામી,
સરવાળો, ગુણાકાર, $\Char{=}$,
પ્રક્ષેપણો અને આમાંથી સંયોજન તથા
નિયમિત વિધેયો પર લાગુ કરેલા
લઘુત્તમીકરણ વડે વ્યાખ્યાયિત
કોઈપણ વિધેય વાપરી શકીએ છીએ.
જે સંબંધનું લાક્ષણિક વિધેય આ
રીતે વ્યાખ્યાયિત થાય તે સંબંધ પણ
વાપરી શકીએ. અગાઉની જેમ,
આ સંબંધો બૂલિયન સંયોજનો અને
સીમિત પરિમાણન હેઠળ બંધ છે તે
બતાવી શકીએ; દાખલા તરીકે:
\begin{align*}
\fn{not}(x) & \defis \Char{=}(x,0)\\
\bmin{x \leq z}{R(x,y)} & \defis \umin{x}{(R(x,y) \lor x = z)}\\
\bexists{x \leq z}{R(x,y)} & \defiff R(\bmin{x \leq z}{R(x,y)}, y)
\end{align*}
પછી આ બધું પણ આદિમ પુનરાવર્તન
વિના વ્યાખ્યાયિત કરી શકાય છે
તે બતાવી શકીએ:
\begin{enumerate}
\item જોડ-નિરૂપણ વિધેય,
  $J(x,y) = \frac{1}{2}[(x+y)(x+y+1)] + x$;
% maybe explain more what is going on here, a bit confusing.
\item પ્રક્ષેપણ વિધેયો
\begin{align*}
K(z) & = \bmin{x \leq z}{\bexists{y \leq z}{z = J(x,y)}},\\
L(z) & = \bmin{y \leq z}{\bexists{x \leq z}{z = J(x,y)}};
\end{align*}
\item ન્યૂનતા સંબંધ $x < y$;
\item વિભાજ્યતા સંબંધ $x \mid y$;
% \item $x \tsub y$
% \item $\fn{Prime}(x)$
% \item Assuming $p$ is prime, the relation ``$x$ is a power of $p$'':
% \[
% \bforall{y \leq x}{(y \mid x \lif y = 1 \lor y = x)}.
% \]
\item $y$ને~$x$ વડે ભાગતાં
  મળતું શેષફળ આપતું વિધેય
  $\fn{rem}(x,y)$.
\end{enumerate}
હવે આ વ્યાખ્યાયિત કરો:
\begin{align*}
\beta^*(d_0,d_1,i) & = \fn{rem}(1+(i+1) d_1,d_0) \text{અને}\\
\beta(d,i) & = \beta^*(K(d),L(d),i).
\end{align*}
આ જ આપણને જોઈતું વિધેય છે.
ઉપરની જેમ $a_0,\dots,a_n$
આપેલાં હોય, તો આ લો:
\[
j = \max(n,a_0+1,\dots,a_n+1),
\]
અને $d_1 = \lcm(1,\dots,j)$ લો.
ઉપરના \olref{rel-prime} પ્રમાણે
$1+d_1$, $1+2 d_1$, \dots,
$1+(n+1) d_1$ જોડી-જોડી
પરસ્પર અવિભાજ્ય છે; અને~\olref{less}
પ્રમાણે તે અનુક્રમે
$a_0,\dots,a_n$ કરતાં મોટાં છે.
સુનઝીના પ્રમેય પ્રમાણે એવું~$d_0$
છે કે દરેક~$i$ માટે
\[
d_0 \equiv a_i \mod (1+(i+1)d_1)
\]
અને તેથી ($d_1$ એ~$a_i$
કરતાં મોટું હોવાથી)
\[
a_i = \fn{rem}(1+(i+1)d_1,d_0).
\]
હવે $d = J(d_0,d_1)$ લો.
તો દરેક $i \le n$ માટે
\begin{align*}
\beta(d,i) & =  \beta^*(d_0,d_1,i) \\
& =  \fn{rem}(1+(i+1) d_1,d_0) \\
& =  a_i
\end{align*}
મળે છે, અને એ જ જરૂરી હતું.
આથી $\beta$-વિધેયના લેમ્માની
સાબિતી પૂરી થાય છે.

\begin{prob}
  સંબંધો $x < y$ અને $x \mid y$
  તથા વિધેય~$\fn{rem}(x,y)$
  આદિમ પુનરાવર્તન વિના
  વ્યાખ્યાયિત કરી શકાય છે તે
  બતાવો. તમે $0$, ઉત્તરગામી,
  સરવાળો, ગુણાકાર, $\Char{=}$,
  પ્રક્ષેપણો તથા સીમિત
  લઘુત્તમીકરણ અને પરિમાણન
  વાપરી શકો છો.
\end{prob}

\end{document}
